\subsection{Описание верифицируемых свойств}

Свойства вынесены в файл \code{properties.pml}, который подключает
модель командой \verb|#include "modbus.pml"|. Всего сформулировано
14 LTL-свойств. Они проверяют обработку корректных и ошибочных запросов,
сохранность памяти, ответы устройств и завершение обмена.

В формулах используются оператор \code{[]} со смыслом «всегда»,
импликация \code{->} и логические связки. Для свойства прогресса
дополнительно используется сильный оператор \code{U}. Формулы записаны
как требования к поведению модели. При генерации верификатора SPIN
автоматически строит по отрицанию выбранной формулы автомат
\code{never claim}, который ищет нарушающее ее выполнение~\cite{ltl}.
Отдельные \code{assert}-выражения и написанный вручную
\code{never claim} в проекте не используются.

Важной особенностью является выбор момента проверки. Условия, описывающие
готовый результат, проверяются при \code{finished}. Ограничения,
которые должны соблюдаться и во время обработки, используют \code{active}.
Это позволяет не сравнивать поля разных транзакций при подготовке нового
запроса. При этом выполнение условий после завершения само по себе еще
не гарантирует, что транзакция завершится: этому посвящено отдельное
свойство прогресса.

\subsubsection{Игнорирование запроса с ошибочной CRC}

Первое свойство требует, чтобы кадр с неверной CRC не приводил ни к записи,
ни к ответу. Для проверки вводится макрос \code{finished}, обозначающий
любое из трех итоговых состояний Master. Макрос \code{no_writes}
проверяет отсутствие операций записи у всех устройств.
\code{unchanged(s)} сравнивает оба регистра одного Slave со снимком,
а \code{all_unchanged} распространяет сравнение на все три устройства.

\listing{defs_p01.pml}{Признаки завершения и сохранности памяти}
\listing{p01_crc_ignored.pml}{Свойство 1 Обработка ошибочной CRC}

Формула проверяется во всех состояниях выполнения. Если обмен завершен
и CRC неверна, число ответов должно быть равно нулю, флаги записи должны
оставаться сброшенными, а память --- совпадать со снимком. Для адреса 0
ожидается завершение широковещательного обмена, для любого другого
адреса --- тайм-аут.

\subsubsection{Обращение к отсутствующему устройству}

Второе свойство относится к адресу 4, для которого процесс Slave не создан.
Такой запрос должен завершаться тайм-аутом, не затрагивая память
существующих устройств и не создавая ответов.

\listing{p02_absent_slave.pml}{Свойство 2 Отсутствующий получатель}

Проверка не ограничивается корректными кадрами: при любом запросе
к отсутствующему адресу результат должен удовлетворять указанным условиям.
Это позволяет отдельно проверить фильтрацию получателя всеми Slave.

\subsubsection{Отсутствие ответа на широковещательный запрос}

Широковещательный запрос не должен создавать ответ даже во время
обработки. Для этого используется \code{active}: он истинен
в состоянии ожидания и в любом итоговом состоянии транзакции.

\listing{defs_p03.pml}{Признак активной транзакции}
\listing{p03_broadcast_silent.pml}{Свойство 3 Молчание при широковещательной передаче}

Условие проверяет одновременно счетчик отправленных ответов и длину
канала. Поэтому нарушение нельзя скрыть последующим извлечением сообщения:
после отправки значение счетчика уже не было бы нулевым.

\subsubsection{Запись корректного широковещательного запроса}

Для проверки записи вводится \code{valid_write}. Он требует верной CRC,
функции \code{FC06} и существующего регистра, но не проверяет адрес
получателя. Макрос \code{written(s)} устанавливает, что нужный регистр
содержит переданное значение, а второй регистр этого Slave не изменился.

\listing{defs_p04.pml}{Корректная запись и ее результат}
\listing{p04_broadcast_written.pml}{Свойство 4 Запись всеми Slave}

После завершения корректного запроса с адресом 0 все три устройства должны
иметь установленный флаг записи и ожидаемое содержимое памяти. Отдельная
проверка флагов существенна, если переданное значение уже находилось
в регистре до запроса. Итоговое состояние Master должно соответствовать
широковещательному завершению.

\subsubsection{Выполнение корректной адресной записи}

Макрос \code{addressed} обозначает адрес существующего устройства,
то есть диапазон от 1 до 3. Пятое свойство требует, чтобы корректная запись
такому устройству завершалась нормальным ответом и действительно
выполнялась у выбранного Slave.

\listing{defs_p05.pml}{Признак существующего адресного получателя}
\listing{p05_addressed_write.pml}{Свойство 5 Успешная адресная запись}

Дизъюнкция из трех ветвей связывает адрес запроса с нужным индексом
в массивах. Для выбранного устройства проверяются факт записи и
\code{written(s)}. Поля нормального ответа проверяются далее отдельной
формулой, а сохранность остальных Slave --- свойством изоляции.

\subsubsection{Изоляция записи}

Адресная операция должна изменять только выбранный регистр выбранного
устройства. Макрос \code{isolated(s)} рассматривает два случая.
Если Slave не является получателем, у него не должно быть ни операции
записи, ни изменения памяти. Если запрос адресован ему, неизменным должен
оставаться другой регистр, а при ошибочном адресе регистра --- оба.

\listing{defs_p06.pml}{Проверка изоляции одного Slave}
\listing{p06_write_isolation.pml}{Свойство 6 Изоляция адресной записи}

Формула применяется ко всем завершенным адресным запросам, включая
ошибочные. Совместная проверка трех экземпляров \code{isolated}
охватывает и посторонние устройства, и невыбранную часть памяти получателя.

\subsubsection{Содержимое нормального ответа}

Седьмое свойство начинается с факта получения нормального ответа.
В таком случае исходный запрос должен быть корректной адресной записью,
а количество отправленных ответов --- равняться одному.

\listing{p07_normal_response.pml}{Свойство 7 Поля нормального ответа}

Проверяется совпадение адреса Slave, адреса регистра и значения с запросом,
а также код функции \code{FC06}. Таким образом, формула исключает
нормальный ответ на поврежденный, широковещательный или неподдерживаемый
запрос и проверяет содержимое уже принятого сообщения.

\subsubsection{Исключение для неподдерживаемой функции}

Для описания исключений используется \code{exception_response(code)}.
Он проверяет состояние получения ответа, единственность отправки,
совпадение адреса и установку старшего бита кода функции. Дополнительно
проверяются нулевые поля регистра и данных и требуемый код исключения.

\listing{defs_p08.pml}{Общий вид ответа с исключением}
\listing{p08_illegal_function.pml}{Свойство 8 Неподдерживаемая функция}

Если адресный запрос имеет верную CRC, но его функция отличается от
\code{FC06}, после завершения требуется исключение 1. Операций записи
быть не должно, все регистры должны совпадать со снимком. В генерируемых
запросах неподдерживаемая функция представлена кодом \code{FC03}.

\subsubsection{Исключение для несуществующего регистра}

Девятое свойство использует тот же макрос ответа, но требует код
исключения 2. Предпосылка включает верную CRC и поддерживаемую функцию,
поэтому проверяется именно ошибка адреса регистра.

\listing{p09_illegal_register.pml}{Свойство 9 Ошибочный адрес регистра}

После такого запроса ответ должен иметь правильное оформление исключения,
а память всех Slave должна остаться неизменной. Флаги записи позволяют
проверить также отсутствие попытки записи без видимого изменения значения.

\subsubsection{Допустимость каждой выполненной записи}

Макрос \code{authorized(s)} связывает установленный флаг записи
с корректностью запроса и его адресом. Если Slave выполнял запись,
запрос должен удовлетворять \code{valid_write} и быть адресован
этому устройству либо всем устройствам сразу.

\listing{defs_p10.pml}{Условие допустимости записи}
\listing{p10_write_authorized.pml}{Свойство 10 Запись только по допустимому запросу}

В отличие от свойств готового результата, это условие проверяется во всех
активных состояниях. Название \code{authorized} здесь означает
допустимость операции по правилам модели: отдельный механизм
аутентификации участников в нее не вводится.

\subsubsection{Единственность ответа}

Один запрос не должен приводить к нескольким ответам. Для проверки
используется накопительный счетчик \code{responses_sent}, который
сбрасывается только при подготовке следующей транзакции.

\listing{p11_single_response.pml}{Свойство 11 Не более одного ответа}

Условие допускает ноль ответов, что необходимо для широковещательной
передачи, неверной CRC и отсутствующего получателя. Одновременно оно
запрещает повторную отправку в любой момент активной транзакции.

\subsubsection{Обязательное завершение начатой транзакции}

Проверок конечного результата недостаточно: они могли бы выполняться
формально и при отсутствии завершенных обменов. Поэтому двенадцатое
свойство требует прогресса из \code{MASTER_WAITING}.

\listing{p12_transaction_progress.pml}{Свойство 12 Прогресс транзакции}

Сильный оператор \code{U} требует, чтобы \code{finished} обязательно
наступил, а до этого Master оставался в состоянии ожидания. Таким образом,
запрос нельзя бесконечно обрабатывать или заменить новым запросом
до фиксации результата.

Проверка выполняется без дополнительного параметра справедливости
планирования \code{-f}. В данной модели каждый Slave проходит конечную
обработку и затем ожидает следующий запрос, а Master ожидает уведомления.
Поэтому бесконечная работа постороннего участника внутри одной транзакции
отдельным циклом не предусмотрена. Полученный результат относится
к этой структуре процессов.

\subsubsection{Согласованность состояния завершения}

Следующее свойство уточняет смысл \code{finished}. К моменту фиксации
результата все обработчики должны завершить запрос, а сообщения текущей
транзакции должны быть извлечены из каналов.

\listing{p13_completion_clean.pml}{Свойство 13 Завершение всех обработчиков и очистка каналов}

Проверяются три признака \code{slave_processed}, три очереди запросов,
канал служебных уведомлений и канал ответов. Это условие не требует
обнуления регистров: полезные данные сохраняются между транзакциями,
а очищается только обмен сообщениями.

\subsubsection{Ошибочный широковещательный запрос}

Последнее свойство объединяет требования к широковещательной передаче
и обработке ошибок. Если адрес равен нулю, но \code{valid_write}
не выполняется, обмен должен завершиться без ответа и изменения памяти.

\listing{p14_invalid_broadcast.pml}{Свойство 14 Ошибочная широковещательная передача}

Формула охватывает все предусмотренные генератором ошибки широковещательного
запроса. Ни один Slave не должен выполнить запись, счетчик и канал ответов
должны остаться пустыми, а Master должен зафиксировать
\code{MASTER_BROADCAST_DONE}. Совместно с четвертым свойством это описывает
оба исхода широковещательной обработки: запись корректного запроса
и отсутствие изменений при ошибке.
